\documentclass{article}

\usepackage{amsmath,amssymb}
\usepackage{xurl}
\usepackage[hidelinks]{hyperref}
\setlength{\emergencystretch}{2em}

% Shared commands for notes in docs/paper-gaps/.

\DeclareUnicodeCharacter{2080}{\ifmmode{}_{0}\else\textsubscript{0}\fi}
\DeclareUnicodeCharacter{2081}{\ifmmode{}_{1}\else\textsubscript{1}\fi}
\DeclareUnicodeCharacter{2082}{\ifmmode{}_{2}\else\textsubscript{2}\fi}
\DeclareUnicodeCharacter{2099}{\ifmmode{}_{n}\else\textsubscript{n}\fi}

\newcommand{\Lean}{\textsc{Lean}}
\ExplSyntaxOn
\NewDocumentCommand{\leanid}{m}{
  \ifmmode
    \textsf{\detokenize{#1}}
  \else
    \group_begin:
    \urlstyle{sf}
    \tl_set:Nn \l_tmpa_tl {#1}
    \tl_replace_all:Nnn \l_tmpa_tl {\_} {_}
    \exp_args:NV \path \l_tmpa_tl
    \group_end:
  \fi
}
\ExplSyntaxOff
\providecommand{\tr}{\operatorname{tr}}
\newcommand{\tnleanIssueBase}{https://github.com/LionSR/TNLean/issues/}
\newcommand{\tnleanPrBase}{https://github.com/LionSR/TNLean/pull/}
\newcommand{\tnleanDocBase}{https://sirui-lu.com/TNLean/docs/}
\newcommand{\ghissue}[1]{\href{\tnleanIssueBase#1}{GitHub issue \##1}}
\newcommand{\ghpr}[1]{\href{\tnleanPrBase#1}{PR \##1}}
\newcommand{\doclink}[1]{\href{\tnleanDocBase#1.html}{\path{#1}}}

% Dirac notation and the identity operator, as in blueprint/src/macros/common.tex.
% \providecommand keeps note-local definitions authoritative.
\usepackage{dsfont}
\providecommand{\ket}[1]{|#1\rangle}
\providecommand{\bra}[1]{\langle #1|}
\providecommand{\braket}[2]{\langle #1 | #2 \rangle}
\providecommand{\ketbra}[2]{|#1\rangle\!\langle #2|}
\providecommand{\Id}{\mathds{1}}
\providecommand{\one}{\mathds{1}}
\providecommand{\C}{\mathbb{C}}
\providecommand{\MN}[1]{M_{#1}(\C)}
\providecommand{\MD}{M_D(\C)}

% External referencing for paper sources.  The build helper writes sanitized
% auxiliary files under build/paper-aux/ and excludes duplicated source labels.
\usepackage{xr}

% Import a generated paper auxiliary file with a caller-supplied label prefix.
% The candidate paths support builds from docs/paper-gaps/, the repository root,
% and build/paper-aux/ itself.
\newcommand{\importpaperaux}[2]{%
  \IfFileExists{../../build/paper-aux/#2.aux}{%
    \externaldocument[#1]{../../build/paper-aux/#2}%
  }{%
    \IfFileExists{build/paper-aux/#2.aux}{%
      \externaldocument[#1]{build/paper-aux/#2}%
    }{%
      \IfFileExists{#2.aux}{%
        \externaldocument[#1]{#2}%
      }{}%
    }%
  }%
}

% General external reference macros
\newcommand{\paperref}[2]{\ref{#1-#2}}
\newcommand{\papereqref}[2]{(\ref{#1-#2})}

% CPSV16 source (1606.00608 / MPDO-22-12-17-2.tex) external document and shortcuts
\importpaperaux{cpsv16-}{1606.00608/MPDO-22-12-17-2-tnlean}
\newcommand{\cpsvref}[1]{\paperref{cpsv16}{#1}}
\newcommand{\cpsveqref}[1]{\papereqref{cpsv16}{#1}}

% Verdict marker: \gapnote{<kind>}{<status>}. Kind is the policy
% classification (clarification, local-correction, scope-restriction,
% unfaithful, false-source, open-gap); status is open, wip (elimination
% underway), resolved, or historical. Machine-read by the site index;
% prints a verdict line.
\newcommand{\gapnote}[2]{\par\noindent\textbf{Verdict:} #1 (#2).\par}


\title{Positive Block Dimensions in the PGVWC07 Canonical Form}
\author{}
\date{2026-09-05}

\begin{document}
\maketitle
\gapnote{local-correction}{resolved}

\section{The printed block form}

Theorem~4 of Ref.~\cite{PerezGarcia2007MPSReps}
(\path{Papers/quant-ph_0608197/MPSarchive.tex}, lines~742--763) writes the
matrices of a translation-invariant canonical representation as
\begin{align}
    A_i
     & =\bigoplus_j\lambda_jA_i^j,
    \qquad 1\ge\lambda_j>0,
    \label{eq:pgvwc07_ti_canonical_form_positive_blocks_block_form}
\end{align}
where every block $A^j$ satisfies $\sum_iA_i^j(A_i^j)^\dagger=\Id$, admits a
diagonal positive full-rank matrix $\Lambda^j$ with
$\sum_i(A_i^j)^\dagger\Lambda^jA_i^j=\Lambda^j$, and has the identity as the
only fixed point of $X\mapsto\sum_iA_i^jX(A_i^j)^\dagger$. The source does not
state the size of a block.

\section{The convention}

A block of size zero contributes nothing to
(\ref{eq:pgvwc07_ti_canonical_form_positive_blocks_block_form}) and is not a
summand the source intends: its weight $\lambda_j$ would be a listed weight of
no matrix, and its conditions would hold vacuously. The formal predicate for
the block form of Theorem~4 therefore requires every listed block to have
positive bond dimension. This is a convention of the same kind as the
nonzero-coefficient convention for canonical forms, not a faithfulness gap:
the predicate admits exactly the block forms the source describes, and the
existence theorem for the canonical form produces only blocks of positive
dimension.

The convention is used once. In the uniqueness theorem for the canonical
form (Theorem~7, lines~1002--1015), condition C1 forces a single block, by the
first assertion of the proposition on condition C1 (lines~911--919): a matrix
unit supported on an off-diagonal block is annihilated by every word. Choosing
that matrix unit requires one bond index in each of two distinct blocks, which
positive block dimension supplies.\footnote{The formal declarations are
    \leanid{MPSTensor.PGVWC07CanonicalFormData} and
    \leanid{MPSTensor.PGVWC07CanonicalFormData.not_isNBlkInjective_of_ne}.}

\bibliographystyle{alpha}
\bibliography{references}

\end{document}
